\section{The inherited physical interface}
\label{sec:tpc36-physical-interface}

We import the actual packet and proved orbit reduction from TPC-25 and
TPC-32--TPC-35 \citep{WangTPC25,WangTPC32,WangTPC33,WangTPC34,WangTPC35}.
This section makes clear which object is reassembled and which later
completion remains open.

\subsection{Rows, coefficients, and the joint multiplier}

Fix \(h\in\Z\setminus\{0\}\), put \(H=\operatorname{rad}|h|\), and work
on one source-compatible dyadic packet with
\begin{equation}
 Q=LD,\qquad QJ\asymp X,\qquad R\le T<U_0\asymp X.
 \label{eq:tpc36-physical-scales}
\end{equation}
A row is
\begin{equation}
 \alpha=(\ell_\alpha,d_\alpha),\qquad
 m_\alpha=\ell_\alpha d_\alpha\asymp Q,
 \label{eq:tpc36-row}
\end{equation}
where \(\ell_\alpha\asymp L\) is prime, \(d_\alpha\asymp D\) is
squarefree, and \((d_\alpha,H)=1\).  The source separation
\begin{equation}
 R=X^{1/2-\delta+o(1)},\qquad
 D\le R^{1/2+o(1)},\qquad
 L>\max\{2R,2|h|\}
 \label{eq:tpc36-source-separation}
\end{equation}
makes the integer \(m_\alpha\) determine its row label uniquely.

The two physical row coefficients are
\begin{equation}
 \gamma_\alpha^{(i)}
 =\mu(d_\alpha)(\log\ell_\alpha)
  \omega_D^{(i)}(d_\alpha)
  \psi_L^{(i)}(\ell_\alpha/L)\zeta_\alpha^{(i)},
 \qquad i=1,2,
 \label{eq:tpc36-row-coefficients}
\end{equation}
with fixed smooth and finite-residue factors.  The inherited masses are
\begin{align}
 \sum_\alpha|\gamma_\alpha^{(i)}|
 &\ll X^{o(1)}Q,
 \label{eq:tpc36-row-l1}\\
 \sum_\alpha|\gamma_\alpha^{(i)}|^2
 &\ll X^{o(1)}Q.
 \label{eq:tpc36-row-l2}
\end{align}
The logarithm in the original \(L^2\) estimate is absorbed into
\(X^{o(1)}\).

For \(j\) in an interval of \(O(J)\) integers, put
\begin{equation}
 N_\alpha(j)=m_\alpha j+h\asymp X.
 \label{eq:tpc36-target}
\end{equation}
On the physical support every divisor \(u\mid N_\alpha(j)\) satisfies
\begin{equation}
 (u,m_\alpha jH)=1.
 \label{eq:tpc36-target-primitivity}
\end{equation}

Let
\begin{equation}
 \Delta=QX^{-\kappa_0},\qquad G=X^{\kappa_0},
 \label{eq:tpc36-pruning-scales}
\end{equation}
where \(\kappa_0>0\) is the inherited pruning parameter.  The actual
static mask is
\begin{equation}
 \boxed{
 \mathfrak m_{\rm act}(\alpha,\gamma)
 =\one_{\ell_\alpha\ne\ell_\gamma}
  \one_{|m_\alpha-m_\gamma|>\Delta}
  \one_{(d_\alpha,d_\gamma)\le G}.}
 \label{eq:tpc36-actual-static-mask}
\end{equation}
The full multiplier is
\begin{equation}
 \mathfrak A_{\alpha,\gamma}^{\rm act}(j)
 =\mathfrak m_{\rm act}(\alpha,\gamma)
  \Xi_{\alpha,\gamma}(j)
  \cW_{\alpha,\gamma}(j/J).
 \label{eq:tpc36-actual-joint-multiplier}
\end{equation}
Let \(H_0=O_h(1)\) be a common period in \(j\) for the fixed-residue
factor \(\Xi\).  TPC-25 proves that \(\Xi\) may be handled by first
splitting the \(O_h(1)\) cells \(j\bmod H_0\), and that on each cell the
smooth factor \(\cW\) has an absolutely convergent Mellin projective
separation with polynomially weighted orbit-seminorm mass bounded
independently of \(X\).  The new issue is therefore the nonsmooth mask
\eqref{eq:tpc36-actual-static-mask}.

\subsection{Orbit slicing and reflected columns}

Let
\begin{equation}
 C_m(j)=\sum_{\substack{T<u\le U_0\\u\mid mj+h}}-\mu(u)\log u.
 \label{eq:tpc36-ultra-increment}
\end{equation}
The left orbit slice is
\begin{equation}
 Y_{\gamma,j}^{L}
 =\sum_\alpha\gamma_\alpha^{(1)}
  \mathfrak A_{\alpha,\gamma}^{\rm act}(j)C_{m_\alpha}(j),
 \qquad
 V_L=\sum_{\gamma,j}|Y_{\gamma,j}^{L}|^2,
 \label{eq:tpc36-left-orbit-slice}
\end{equation}
and the right slice exchanges the two row systems.  TPC-34 proves that
the bound \eqref{eq:tpc36-inherited-gate-intro} is sufficient for the
remaining scalar shell and proves the diagonal
\eqref{eq:tpc36-inherited-diagonal-intro}.

Open the ultra divisor in \eqref{eq:tpc36-left-orbit-slice}.  For a
nonzero term write
\begin{equation}
 \ell dj+h=su,
 \qquad
 T<u\le U_0,\qquad
 1\le s\ll X/T.
 \label{eq:tpc36-reflection}
\end{equation}
After one fixes \(\gamma\), a product-ready row decomposition would turn
the inner contribution into layers of the schematic form
\begin{equation}
 \cZ_{\gamma,\rho}(j,s)
 =\sum_{\ell,d,u}
 b_{\gamma,\rho}(\ell)f_{\gamma,\rho}(d)
 g_{\gamma,\rho}(u)
 \one_{su=\ell jd+h}.
 \label{eq:tpc36-integer-product-slope}
\end{equation}
The goal of \cref{sec:tpc36-fixed-column} is to prove the exact
separation in \(\ell,d,j\) that precedes
\eqref{eq:tpc36-integer-product-slope}.

\subsection{What static reassembly does not include}

Modulo a prime \(q\nmid h\), a complete version of
\eqref{eq:tpc36-integer-product-slope} becomes
\begin{align}
 R_s(a)&=\sum_{D\in\F_q}f(D)
 g\!\left(s^{-1}(aD+h)\right),
 \label{eq:tpc36-Rsa-interface}\\
 (Tb)(j,s)&=\sum_{\ell\in\F_q^\times}b(\ell)R_s(\ell j).
 \label{eq:tpc36-T-interface}
\end{align}
TPC-35 gives the exact character spectrum of
\eqref{eq:tpc36-T-interface}.  Passing from
\eqref{eq:tpc36-integer-product-slope} to that complete finite-field
operator additionally requires:
\begin{enumerate}[label=(\roman*)]
 \item controlling aliases between the integer equation and its modular
 relaxation;
 \item completing the physical \(j,d,u,s\) ranges without a polynomial
 cost; and
 \item retaining the coherent \(s\)-sum in the topology of
 \eqref{eq:tpc36-left-orbit-slice}.
\end{enumerate}
None of these statements follows merely from a low projective norm for
\eqref{eq:tpc36-actual-static-mask}.  We keep them separate throughout.

\subsection{Endpoint scales}

At the high--\(\beta\) endpoint,
\begin{equation}
 L=X^{99979/210000+o(1)},\qquad
 D=X^{10049/52500+o(1)},
 \label{eq:tpc36-LD-endpoint}
\end{equation}
in addition to \eqref{eq:tpc36-endpoint-intro}.  Hence
\begin{equation}
 D=o(J),\qquad J=o(Q),\qquad QJ=X^{1+o(1)}.
 \label{eq:tpc36-endpoint-relations}
\end{equation}
The number of formal column layers \(O(D)\) below is polynomial, but its
exponent is irrelevant to estimates that reassemble by signed projective
mass.  Counting layers absolutely would be the wrong norm.
